%%%%%%%%%%%%%%%%%%%%%%%%%%%%%%%%%%%%%%%%%%%%%%%%%%%%%%%%%%%%%%%%%%%%%%%%%%%%%%%%%%
\begin{frame}[fragile]\frametitle{}
\begin{center}
{\Large Decorators}
\end{center}
\end{frame}

%%%%%%%%%%%%%%%%%%%%%%%%%%%%%%%%%%%%%%%%%%%%%%%%%%%%%%%%%%%%%%%%%%%%%%%%%%%%%%%%%%%
\begin{frame}[fragile]\frametitle{Decorators Intro}
``This amazing feature appeared in the language almost apologetically and with concern that it might not be that useful.''
-  Bruce Eckel in ``An Introduction to Python Decorators''
\end{frame}

%%%%%%%%%%%%%%%%%%%%%%%%%%%%%%%%%%%%%%%%%%%%%%%%%%%%%%%%%%%%%%%%%%%%%%%%%%%%%%%%%%%
\begin{frame}[fragile]\frametitle{Decorators Intro}
        \begin{lstlisting}
def my_decorator(some_function):
    def wrapper():
        print("Something is happening before some_function() is called.")
        some_function()
        print("Something is happening after some_function() is called.")
    return wrapper

def just_some_function():
    print("Wheee!")

just_some_function = my_decorator(just_some_function)
just_some_function()

Something is happening before some_function() is called.
Wheee!
Something is happening after some_function() is called.
\end{lstlisting}
\end{frame}



%%%%%%%%%%%%%%%%%%%%%%%%%%%%%%%%%%%%%%%%%%%%%%%%%%%%%%%%%%%%%%%%%%%%%%%%%%%%%%%%%%%
\begin{frame}[fragile]\frametitle{Decorators Intro}
    \begin{itemize}
    	\item Decorator is way to dynamically add some new behavior to some objects. 
    \item  Since functions and classes are objects, they can be passed around. 
    \item Since they are mutable objects, they can be modified. 
    \item The act of altering a function or class object after it has been constructed but before it is bound to its name is called decorating.
    \end{itemize}
\end{frame}

%%%%%%%%%%%%%%%%%%%%%%%%%%%%%%%%%%%%%%%%%%%%%%%%%%%%%%%%%%%%%%%%%%%%%%%%%%%%%%%%%%%
\begin{frame}[fragile]\frametitle{Intuition: Decorators}
\begin{block}{Intuition}
A decorator is like gift-wrapping a present: the original gift (the function) is still inside,
unchanged, but it's now wrapped in extra paper (extra behavior) that runs before -- and possibly
after -- you actually get to the gift. The \texttt{@decorator} syntax is just shorthand for
handing your function to the wrapping function and using whatever comes back in its place.
\end{block}
\end{frame}

%%%%%%%%%%%%%%%%%%%%%%%%%%%%%%%%%%%%%%%%%%%%%%%%%%%%%%%%%%%%%%%%%%%%%%%%%%%%%%%%%%%
\begin{frame}[fragile]\frametitle{Decorators Intro}
        \begin{lstlisting}
def my_decorator(some_function):
    def wrapper():
        num = 10
        if num == 10:
            print("Yes!")
        else:
            print("No!")
        some_function()
        print("Something is happening after some_function() is called.")
    return wrapper

def just_some_function():
    print("Wheee!")

just_some_function = my_decorator(just_some_function)
just_some_function()

Yes!
Wheee!
Something is happening after some_function() is called.
\end{lstlisting}
\end{frame}

%%%%%%%%%%%%%%%%%%%%%%%%%%%%%%%%%%%%%%%%%%%%%%%%%%%%%%%%%%%%%%%%%%%%%%%%%%%%%%%%%%%
\begin{frame}[fragile]\frametitle{Decorators Intro}
Python allows you to simplify the calling of decorators using the @ symbol (this is called “pie” syntax).
        \begin{lstlisting}
@my_decorator
def just_some_function():
    print("Wheee!")

just_some_function()

Yes!
Wheee!
Something is happening after some_function() is called.
\end{lstlisting}
\end{frame}

%%%%%%%%%%%%%%%%%%%%%%%%%%%%%%%%%%%%%%%%%%%%%%%%%%%%%%%%%%%%%%%%%%%%%%%%%%%%%%%%%%%
\begin{frame}[fragile]\frametitle{Decorators Intro}
There are two things hiding behind the name ``decorator'': the function which does the work of
decorating -- performs the real work -- (marked \texttt{[2]} below), and the expression
adhering to decorator syntax -- an at-symbol and the decorating function's name -- applied just
above a normally-defined function (marked \texttt{[1]}):
        \begin{lstlisting}
@decorator             # [2]
def function():        # [1]
    pass
\end{lstlisting}
\end{frame}

%%%%%%%%%%%%%%%%%%%%%%%%%%%%%%%%%%%%%%%%%%%%%%%%%%%%%%%%%%%%%%%%%%%%%%%%%%%%%%%%%%%
\begin{frame}[fragile]\frametitle{Decorators Intro}
    \begin{itemize}
    \item The part after @ must be a simple expression, usually this is just the name of a function or class. 
    \item This part is evaluated first, and after the function defined below is ready,
    \item The decorator is called with the newly defined function object as the single argument. 
    \item The value returned by the decorator is attached to the original name of the function.
    \end{itemize}
\end{frame}

%%%%%%%%%%%%%%%%%%%%%%%%%%%%%%%%%%%%%%%%%%%%%%%%%%%%%%%%%%%%%%%%%%%%%%%%%%%%%%%%%%%
\begin{frame}[fragile]\frametitle{Decorators Intro}
    \begin{itemize}
    \item Decorators can be applied to functions and to classes. 
    \item For classes the semantics are identical : the original class definition is used as an argument to call the decorator and 
    \item whatever is returned is assigned under the original name.
    \end{itemize}
\end{frame}

%%%%%%%%%%%%%%%%%%%%%%%%%%%%%%%%%%%%%%%%%%%%%%%%%%%%%%%%%%%%%%%%%%%%%%%%%%%%%%%%%%%
\begin{frame}[fragile]\frametitle{Decorators Example}
In the example we will create a simple example which will print some statement before and after the execution of a function.
        \begin{lstlisting}
>>> def my_decorator(func):
...     def wrapper(*args, **kwargs):
...         print("Before call")
...         result = func(*args, **kwargs)
...         print("After call")
...         return result
...     return wrapper
...
>>> @my_decorator
... def add(a, b):
...     "Our add function"
...     return a + b
...
>>> add(1, 3)
Before call
After call
4
\end{lstlisting}
\end{frame}



%%%%%%%%%%%%%%%%%%%%%%%%%%%%%%%%%%%%%%%%%%%%%%%%%%%%%%%%%%%%%%%%%%%%%%%%%%%%%%%%%%%
\begin{frame}[fragile]\frametitle{Decorators Example}
        \begin{lstlisting}
import time
def timing_function(some_function):
    def wrapper():
        t1 = time.time()
        some_function()
        t2 = time.time()
        return "Time it took to run the function: " + str((t2 - t1)) + "\n"
    return wrapper

@timing_function
def my_function():
    num_list = []
    for num in (range(0, 10000)):
        num_list.append(num)
    print("\nSum of all the numbers: " + str((sum(num_list))))

print(my_function())
\end{lstlisting}
\end{frame}



%%%%%%%%%%%%%%%%%%%%%%%%%%%%%%%%%%%%%%%%%%%%%%%%%%%%%%%%%%%%%%%%%%%%%%%%%%%%%%%%%%%
\begin{frame}[fragile]\frametitle{Decorator Stacking}
    \begin{itemize}
    \item Decorators can be \textbf{stacked}: applied bottom-to-top (inside-out).
    \item The originally defined function is passed to the innermost decorator; each outer decorator receives the result of the one below.
    \item The final return value is bound to the original function name.
    \end{itemize}
\begin{lstlisting}
@decorator_a
@decorator_b
def func(): pass
# Equivalent to: func = decorator_a(decorator_b(func))
\end{lstlisting}
\end{frame}

%%%%%%%%%%%%%%%%%%%%%%%%%%%%%%%%%%%%%%%%%%%%%%%%%%%%%%%%%%%%%%%%%%%%%%%%%%%%%%%%%%%
\begin{frame}[fragile]\frametitle{Decorator Readability}
    \begin{itemize}
    \item The \texttt{@} syntax was chosen for readability.
    \item Placed \emph{before} the function header, it clearly signals that the function is being wrapped.
    \item Multiple decorators are each on their own line, making the chain easy to read.
    \item The only requirement: the decorator must be callable with one argument (the decorated object).
    \end{itemize}
\end{frame}

%%%%%%%%%%%%%%%%%%%%%%%%%%%%%%%%%%%%%%%%%%%%%%%%%%%%%%%%%%%%%%%%%%%%%%%%%%%%%%%%%%%
\begin{frame}[fragile]\frametitle{Decorators Intro}
    \begin{itemize}
    \item The only requirement on decorators is that they can be called with a single argument. 
    \item This means that decorators can be implemented as normal functions, or as classes with a \_\_call\_\_ method, or in theory, even as lambda functions.
    \end{itemize}
\end{frame}

%%%%%%%%%%%%%%%%%%%%%%%%%%%%%%%%%%%%%%%%%%%%%%%%%%%%%%%%%%%%%%%%%%%%%%%%%%%%%%%%%%%
\begin{frame}[fragile]\frametitle{Decorators implemented as classes and as functions}
    \begin{itemize}
    \item The decorator expression (the part after @) can be either just a name, or a call. 
    \item The bare-name approach is nice (less to type, looks cleaner, etc.), but is only possible when no arguments are needed to customise the decorator. 
    \end{itemize}
            \begin{lstlisting}
def simple_decorator(function):
  print("doing decoration")
  return function
  
@simple_decorator
def function():
  print("inside function")

doing decoration
>>> function()
inside function
\end{lstlisting}
\end{frame}

%%%%%%%%%%%%%%%%%%%%%%%%%%%%%%%%%%%%%%%%%%%%%%%%%%%%%%%%%%%%%%%%%%%%%%%%%%%%%%%%%%%
\begin{frame}[fragile]\frametitle{Decorators implemented as classes and as functions}
            \begin{lstlisting}
def decorator_with_arguments(arg):
  print("defining the decorator")
  def _decorator(function):
      # in this inner function, arg is available too
      print("doing decoration, %r" % arg)
      return function
  return _decorator
@decorator_with_arguments("abc")
def function():
  print("inside function")

defining the decorator
doing decoration, 'abc'
function()
inside function
\end{lstlisting}
\end{frame}

%%%%%%%%%%%%%%%%%%%%%%%%%%%%%%%%%%%%%%%%%%%%%%%%%%%%%%%%%%%%%%%%%%%%%%%%%%%%%%%%%%%
\begin{frame}[fragile]\frametitle{Decorators Example}
One of the most used decorators in Python is the login\_required() decorator, which ensures that a user is logged in/properly authenticated before s/he can access a specific route (/secret, in this case):
        \begin{lstlisting}
from functools import wraps
from flask import g, request, redirect, url_for

def login_required(f):
    @wraps(f)
    def decorated_function(*args, **kwargs):
        if g.user is None:
            return redirect(url_for('login', next=request.url))
        return f(*args, **kwargs)
    return decorated_function

@app.route('/secret')
@login_required
def secret():
    pass
\end{lstlisting}
\end{frame}


%%%%%%%%%%%%%%%%%%%%%%%%%%%%%%%%%%%%%%%%%%%%%%%%%%%%%%%%%%%%%%%%%%%%%%%%%%%%%%%%%%%
\begin{frame}[fragile]\frametitle{Decorators Example}
Did you notice that the function gets passed to the functools.wraps() decorator? This simply
preserves the metadata of the wrapped function. Let's look at one last use case -- the following
Flask route handler ensures that the key \texttt{student\_id} is part of the request. Although
this validation works, it really does not belong in the function itself -- and perhaps there are
other routes that need the exact same validation:
        \begin{lstlisting}
@app.route('/grade', methods=['POST'])
def update_grade():
    json_data = request.get_json()
    if 'student_id' not in json_data:
        abort(400)
    # update database
    return "success!"
\end{lstlisting}
\end{frame}

%%%%%%%%%%%%%%%%%%%%%%%%%%%%%%%%%%%%%%%%%%%%%%%%%%%%%%%%%%%%%%%%%%%%%%%%%%%%%%%%%%%
\begin{frame}[fragile]\frametitle{Quick Check: Decorators}
\begin{block}{Think About It}
A colleague writes \lstinline|@login_required| above three different view functions, then
notices all three views seem to share the exact same \texttt{wrapper} name in stack traces. Why,
and is that a problem?
\end{block}
\pause
\begin{block}{Answer}
Every call to \lstinline|login_required(f)| defines and returns a brand-new local function
object, so each decorated view actually gets its own distinct function -- the shared \emph{name}
is just cosmetic. It becomes a real problem for debugging and introspection (stack traces,
\lstinline|help()|, frameworks that key routing off \texttt{\_\_name\_\_}) unless the decorator
uses \lstinline|@functools.wraps(f)| to copy over the original function's \texttt{\_\_name\_\_},
\texttt{\_\_doc\_\_}, etc. -- exactly what the \lstinline|login_required| example does.
\end{block}
\end{frame}

